\section{信頼区間}
\subsection{授業内容}
\subsubsection{科目の中でのこのコマの位置づけ}
標本統計量が確率変数であるということは，心理学実験や調査などで得られるデータとその統計量をそのまま一般化できないことにつながる。
正規母集団から得られた標本平均は幸いにして，母平均の推定量として好ましい性質を有しているが，それにしても標本平均すなわち母平均と考えるのは極端であると言わざるを得ない。
標本統計量が確率変数であり，確率分布(標本分布)の性質がわかっているのなら，それを使って区間推定することが考えられる。
ここでは区間推定についてもシミュレーションを通じて理解する。
\subsubsection{コマ主題細目}
\begin{description}
	\item[点推定と区間推定] 点推定と区間推定の区別を確認した上で，シミュレーションによって点推定が母数と一致する割合，区間推定が母数を含む割合を確認する。前者がゼロであることから，区間推定の必要性が理解でき，かつどのように区間を設定すれば良いかという問題に気づく。
		\begin{flushright}
			$\to$\textcite{kosugi2023}のP.145--146
		\end{flushright}
	\item[母分散が既知の場合] すでに「データ解析基礎」において標準正規分布を利用した区間推定について，95\%区間に対応する$\pm 1.96$という数字を使った例を学んだ。今回はシミュレーションを通じて，標本平均$\pm 1.96\frac{\sigma}{\sqrt{n}}$が正しく母平均を含んでいる回数をチェックし，このことを確認する。
	\item[母分散が未知の場合] 同じく，母分散が未知の場合はt分布を使って区間推定をするのであった。ここでもシミュレーションを通じて，区間推定が正しく母平均を含んでいる回数をチェックし，このことを確認する。可視化することでこのことが理解を深める。
		\begin{flushright}
			$\to$\textcite{kosugi2023}のP.146--156.
		\end{flushright}
\end{description}
\subsubsection{キーワード}
\begin{itemize}
	\item 点推定
	\item 区間推定
	\item t分布
\end{itemize}
\subsection{授業情報}
\paragraph{コマの展開方法}
講義/遠隔/演習
\subsubsection{予習・復習課題}
\paragraph{予習} 一年時の「データ解析基礎」のテキストを復習しておくこと。第17講(後期第二回目)が本時に対応する。
\paragraph{復習} t分布の数式について，テキストのコラムを一読しておく。余裕があれば，テキスト4章後半の相関係数の信頼区間などにも目を向けるとよい。